% !PRINTED 656
\weberpage{656}

Nehmen wir $n = p$ an, so wird $\mathfrak{p}_{p\xi} = \mathfrak{p}_\xi$ (\S.\,172), also ist auch $k_{p\xi} = k_\xi$, und aus (7) und (8) folgt:
\weq{9}{k(p-1) \equiv 0 \ (\mathrm{mod}\ m).}

Hieraus ergeben sich nun wichtige Folgerungen: Wenn zunächst $p - 1$ nicht durch $q$ theilbar ist, so folgt, dass $k$ durch $m$ theilbar sein muss, und wir haben also den Satz:

4. \gesperrt{Die Primfactoren einer Primzahl $p$, die nach dem Modul $q$ nicht mit 1 congruent ist, sind in $\omega$ nur in solchen Potenzen enthalten, deren Exponenten durch $m$ theilbar sind.}

Wenn in dem Falle, wo $m$ eine Potenz von 2 ist, $p - 1$ durch 2, aber nicht durch 4 theilbar ist, wenn also $p \equiv -1 \ (\mathrm{mod}\ 4)$ ist, so folgt aus (9) nur, dass $k$ durch $\tfrac{1}{2} m$, nicht aber, dass $k$ durch $m$ theilbar ist. Der Exponent der in $\omega$ enthaltenen Potenz von $\mathfrak{p}_\xi$ ist, von Vielfachen von $m$ abgesehen, gleich $k$, und wenn $k$ durch $m$ theilbar ist, so ist alles wie im Falle 4. Ist aber $k$ nur durch $\tfrac{1}{2} m$, nicht durch $m$ theilbar, also
\[
k \equiv \tfrac{1}{2} m \ (\mathrm{mod}\ m),
\]
so ist $k - \tfrac{1}{2} m$ durch $m$ theilbar, und wir können mit Rücksicht auf die Zerlegung (3) den Satz so formuliren:

5. \gesperrt{Ist $m$ eine Potenz von 2 und $p$ eine Primzahl congruent mit $-1$ nach dem Modul 4, so lässt sich der Exponent $h$ so bestimmen, dass alle Primfactoren von $p$ in}
\[
\omega \sqrt{p}^{-hm}
\]
\gesperrt{nur zu solchen Potenzen enthalten sind, deren Exponent durch $m$ theilbar ist.}

Es sei endlich $p - 1$ durch $q$, oder, falls $m$ gerade ist, durch 4 theilbar. Ist $m_1$ der grösste gemeinschaftliche Theiler von $m$ und $p - 1$, und ist
\[
m = m_1 m_2,
\]
so muss $k$ nach (9) durch $m_2$ theilbar sein, und wenn wir
\[
k = h m_2
\]
setzen, so ist in $\omega$ nach (8) das Product enthalten:
\weq{10}{\left(\prod^\xi \mathfrak{p}_\xi^{\xi'}\right)^{h m_2}.}
Ausserdem kommen in $\omega$ die Primfactoren $\mathfrak{p}_\xi$ nur noch in solchen Potenzen vor, deren Exponenten durch $m$ theilbar sind.
